% problem_id: S046-thm:freudenthal
% source_id: S046 (HoTT Book)
% source_file: homotopy.tex
% statement_locator: homotopy.tex lines 1642--1645
% proof_locator: homotopy.tex lines 1676--1707
% env: thm
% pairing: tight (verified: the next environment after the statement is its proof)
% status: complete (statement + author proof verbatim + reason + locators)
%
\begin{thm}[The Freudenthal suspension theorem]\label{thm:freudenthal}
  Suppose that $X$ is $n$-connected and pointed, with $n\geq 0$.
  Then the map $\sigma:X\to \Omega\susp(X)$ is $2n$-connected.
\end{thm}

\begin{proof}[Proof of \cref{thm:freudcode}]
  We define $\code(y,p)$ by induction on $y:\susp X$, where the first two cases are~\eqref{eq:freudcodeN} and~\eqref{eq:freudcodeS}.
  It remains to construct, for each $x_1:X$, a dependent path
  \[ \dpath{\lam{y}(\north=y)\to\type}{\merid(x_1)}{\code(\north)}{\code(\south)}. \]
  By \cref{thm:dpath-arrow}, this is equivalent to giving a family of paths
  \[ \prd{q:\north=\south} \code(\north)(\transfib{\lam{y}(\north=y)}{\opp{\merid(x_1)}}{q}) = \code(\south)(q). \]
  And by univalence and transport in path types, this is equivalent to a family of equivalences
  \[ \prd{q:\north=\south} \code(\north,q \ct \opp{\merid(x_1)}) \eqvsym \code(\south,q). \]
  We will define a family of maps
  \begin{equation}\label{eq:freudmap}
    \prd{q:\north=\south} \code(\north,q \ct \opp{\merid(x_1)}) \to \code(\south,q).
  \end{equation}
  and then show that they are all equivalences.
  Thus, let $q:\north=\south$; by the universal property of truncation and the definitions of $\code(\north,\blank)$ and $\code(\south,\blank)$, it will suffice to define for each $x_2:X$, a map
  \begin{equation*}
    \big(\merid(x_2)\ct \opp{\merid(x_0)} = q \ct \opp{\merid(x_1)}\big)
    \to \trunc{2n}{\tsm{x:X} (\merid(x) = q)}.
  \end{equation*}
  Now for each $x_1,x_2:X$, this type is $2n$-truncated, while $X$ is $n$-connected.
  Thus, by \cref{thm:wedge-connectivity}, it suffices to define this map when $x_1$ is $x_0$, when $x_2$ is $x_0$, and check that they agree when both are $x_0$.

  When $x_1$ is $x_0$, the hypothesis is $r:\merid(x_2)\ct \opp{\merid(x_0)} = q \ct \opp{\merid(x_0)}$.
  Thus, by canceling $\opp{\merid(x_0)}$ from $r$ to get $r':\merid(x_2)=q$, so we can define the image to be $\tproj{2n}{(x_2,r')}$.

  When $x_2$ is $x_0$, the hypothesis is $r:\merid(x_0)\ct \opp{\merid(x_0)} = q \ct \opp{\merid(x_1)}$.
  Rearranging this, we obtain $r'':\merid(x_1)=q$, and we can define the image to be $\tproj{2n}{(x_1,r'')}$.

  Finally, when both $x_1$ and $x_2$ are $x_0$, it suffices to show the resulting $r'$ and $r''$ agree; this is an easy lemma about path composition.
  This completes the definition of~\eqref{eq:freudmap}.
  To show that it is a family of equivalences, since being an equivalence is a mere proposition and $x_0:\unit\to X$ is (at least) $(-1)$-connected, it suffices to assume $x_1$ is $x_0$.
  In this case, inspecting the above construction we see that it is essentially the $2n$-truncation of the function that cancels $\opp{\merid(x_0)}$, which is an equivalence.
\end{proof}
